% TOP_PROBLEM: {"id":30007527,"problem_number":"LOCAL-30007527","title":"Quantitative tightening transmission","category_id":21,"category":{"id":21,"name":"economics","display_name":"Economics","description":"Open economic theory statements, published research agendas, and proposed scoped specifications; question provenance and historical priority are qualified per record.","slug":"economics","order_index":21,"created_at":"2026-10-08T00:00:00Z"},"status":"research_question","external_url":null,"legacy_ids":[],"published":false,"statement":"Identify whether matched easing and tightening surprises have opposite firm-level effects at a fixed scale, maturity segment, and liquidity state.","economics":{"original_id":16,"field":"Inflation and monetary policy","kind":"identification","formalization_basis":"adapted_research_question","source_status":"research_agenda","priority_status":"not_established","priority_note":"Original priority is not established. No dated primary statement of this precise problem was verified.","earliest_verified_reference":{"authors":["European Central Bank"],"title":"Research agenda: Monetary policy and its implementation","year":null,"url":"https://www.ecb.europa.eu/pub/economic-research/research_agenda/monetary_policy/html/index.en.html","doi":null,"locator":"“Policy normalisation,” selected areas of focus","evidence_summary":"The agenda asks about the desirable pace and combination of policy-rate and balance-sheet normalisation and implications of banking conditions. QE–QT symmetry is an adaptation, not a verbatim listed problem."},"current_reference":{"authors":["Egemen Eren","Denis Gorea","Daojing Zhai"],"title":"How do quantitative easing and tightening affect firms?","year":2025,"url":"https://www.bis.org/publications/working-paper-1286-how-do-quantitative-easing-and-tightening-affect-firms","doi":null,"locator":"“Focus,” “Contribution,” “Findings,” and abstract on the primary BIS Working Paper 1286 page","evidence_summary":"The paper identifies maturity-specific balance-sheet surprises and estimates firm financing and real responses over 2011–2024. It gives conditional evidence rather than an original symmetry conjecture."},"formalization_note":"Symmetry is a scoped adaptation of the normalisation agenda. The current paper provides relevant firm-level estimates and is not cited as an original open statement."},"statusQualification":"Published question or research agenda verified at the cited locator. The mathematical specification is adapted for this catalogue and includes additional assumptions; source publication does not establish first-ever priority or certify this exact model remains unsolved. Symmetry is a scoped adaptation of the normalisation agenda. The current paper provides relevant firm-level estimates and is not cited as an original open statement.","statusReviewedAt":"2026-10-08"}
% FIRST_POSTED: 2026-10-08
% ENTRY_KIND: statement_only
% !TeX program = lualatex
\documentclass[11pt]{article}
\usepackage{fontspec}
\usepackage[a4paper,margin=25mm]{geometry}
\usepackage{amsmath,amssymb,mathtools}
\usepackage{xurl}
\usepackage[hidelinks]{hyperref}
\newcommand{\sref}[2]{\hyperlink{source:#1:#2}{[S#2]}}
\setlength{\emergencystretch}{3em}
\begin{document}
% problemId:30007527
\section*{Quantitative tightening transmission}\label{30007527}
\subsection{Definitions and mathematical statement}
Let \(b_t\) be an unanticipated central-bank purchase of duration-adjusted government securities, with purchases positive and sales or runoff surprises negative. Fix a maturity segment and define predetermined state \(s_t\) to include reserve abundance, intermediary balance-sheet capacity, and market stress. Let \(F_{j,t+h}(b)\) be firm \(j\)'s potential financing cost, debt maturity, or investment at horizon \(h\), with other policy surprises held fixed. For \(a>0\), define
\[
A_h(a,s)=E[F_{j,t+h}(a)-F_{j,t+h}(0)\mid s_t=s]
+E[F_{j,t+h}(-a)-F_{j,t+h}(0)\mid s_t=s].
\]
Quantitative easing and tightening are symmetric at this scale and state if \(A_h(a,s)=0\). Identify the function \(A_h\), with confidence regions, from announcement surprises that separate duration supply, anticipated runoff, and policy-rate information. Match or reweight the support of easing and tightening episodes; if states do not overlap, report bounds or model-based extrapolation rather than an unconditional symmetry conclusion. Distinguish financing responses from real activity and report heterogeneity by predetermined credit quality. The target is scale- and state-specific symmetry, including possible liquidity thresholds, not a claim that every purchase and sale should mechanically have equal and opposite effects.

\par\textbf{Formulation and scope.} Symmetry is a scoped adaptation of the normalisation agenda. The current paper provides relevant firm-level estimates and is not cited as an original open statement.

\par\textbf{Open-question provenance.} An explicit question or research agenda was verified at the source locator. This does not by itself certify that every later version remains unresolved \sref{30007527}{1}.

\par\textbf{Historical priority.} Original priority is not established. No dated primary statement of this precise problem was verified.

\subsection{Short English statement}
Identify whether matched easing and tightening surprises have opposite firm-level effects at a fixed scale, maturity segment, and liquidity state.

\subsection{Sources}
\begin{itemize}\raggedright
\item[\textnormal{[S1]}]\hypertarget{source:30007527:1}{} European Central Bank. Year not verified. Research agenda: Monetary policy and its implementation. “Policy normalisation,” selected areas of focus.
\url{https://www.ecb.europa.eu/pub/economic-research/research_agenda/monetary_policy/html/index.en.html}.
{\small\textit{Earliest verified question/agenda reference; historical priority qualified below. The agenda asks about the desirable pace and combination of policy-rate and balance-sheet normalisation and implications of banking conditions. QE–QT symmetry is an adaptation, not a verbatim listed problem.}}

\item[\textnormal{[S2]}]\hypertarget{source:30007527:2}{} Egemen Eren, Denis Gorea, Daojing Zhai. 2025. How do quantitative easing and tightening affect firms?. “Focus,” “Contribution,” “Findings,” and abstract on the primary BIS Working Paper 1286 page.
\url{https://www.bis.org/publications/working-paper-1286-how-do-quantitative-easing-and-tightening-affect-firms}.
{\small\textit{Supporting or current literature; see evidence qualification. The paper identifies maturity-specific balance-sheet surprises and estimates firm financing and real responses over 2011–2024. It gives conditional evidence rather than an original symmetry conjecture.}}
\end{itemize}
\end{document}
